\section{Hessian mass and global comparison}
\label{sec:hessian-mass-comparison}

The maximum inequality compares measures on complementary Borel sets.
Integrating it leaves a difference of total masses. The recursive
currents express this difference through the complex Hessians of the
background powers, extending the mass calculation of
\cref{chap:mass} to bounded admissible potentials.

For bounded admissible $t_1,\ldots,t_m$, put
\[
 \mathcal M_m(t_1,\ldots,t_m)=\int_XQ_m(t_1,\ldots,t_m),
 \qquad H_m(t)=Q_m(t,\ldots,t).
\]
Here $H_m(t)$ is a finite positive Radon measure and
$\mathcal M_m$ is the total mass of a mixed positive measure, by
\cref{hessian:bounded-shifted-positive}. For a nonempty subset
$I=\{i_1<\cdots<i_q\}\subset\{1,\ldots,m\}$, put
$i(I)=i_q$ and $I^-=I\setminus\{i_q\}$. Use $D_{I^-}$ for the
recursive current with those ordered inputs, with $D_\varnothing=1$.

\Needspace{17\baselineskip}
\begin{lemma}\label{hessian:bounded-mass-expansion}\index{Hessian measure!total mass}
For every bounded admissible list,
\begin{equation}\label{eq:hessian-bounded-mass-expansion}
 \mathcal M_m(t_1,\ldots,t_m)
 =\int_X\omega^n+
 \sum_{\varnothing\ne I\subset\{1,\ldots,m\}}
 \int_X t_{i(I)}D_{I^-}\wedge\ddc(\omega^{n-|I|}).
\end{equation}
Every integral on the right is an integral of bounded canonical
values against signed measure coefficients.
\end{lemma}
\begin{proof}
Expand $Q_m$ using \eqref{eq:hessian-bounded-shifted-current}.
The term for $I$ pairs $D_I$ with the smooth form
$\omega^{n-|I|}$. By \eqref{eq:hessian-bounded-recursion},
\[
 \langle D_I,\omega^{n-|I|}\rangle
 =\langle t_{i(I)}D_{I^-},\ddc(\omega^{n-|I|})\rangle.
\]
Since $|I^-|\leq m-1$, \cref{hessian:bounded-recursive-currents}
gives measure coefficients for $D_{I^-}$; their bounded product
with $t_{i(I)}$ is well defined. Thus the displayed pairing is
exactly the integral in the sum, and is finite on compact $X$.
The empty subset contributes $\int_X\omega^n$.
\end{proof}

Let $\mu_r$ be the repeated-slot mixed measure from
\cref{hessian:bounded-maximum-measure}, for fixed $r$ and
$z_1,\ldots,z_{m-r}$.
The total-mass difference below therefore has the explicit expression
\eqref{eq:hessian-bounded-mass-expansion} for the two indicated lists.

\Needspace{11\baselineskip}
\begin{proposition}\label{hessian:comparison-mass-defect}\index{Comparison principle!Hessian mass defect}
If $u,v$ are bounded admissible and $w=\max\{u,v\}$, then
\[
 \int_{\{u<v\}}\mathrm d\mu_r(v)
 \leq\int_{\{u<v\}}\mathrm d\mu_r(u)
                +\int_X\mathrm d\mu_r(w)-\int_X\mathrm d\mu_r(u).
\]
\end{proposition}
\begin{proof}
Set $E=\{u<v\}$. The contact-set version of
\cref{hessian:bounded-maximum-measure} gives
\[
 \mu_r(w)\geq\mathbf1_{X\setminus E}\mu_r(u)
                     +\mathbf1_E\mu_r(v).
\]
All these measures are finite and positive. Integrating on $X$ and
using $\mu_r(u)(X\setminus E)=\mu_r(u)(X)-\mu_r(u)(E)$ gives
exactly the assertion. No boundary or closed-background hypothesis
enters this calculation.
\end{proof}

The unmodified global comparison principle holds precisely when the
powers occurring in this mass formula are $\ddc$-closed. Only the
powers $n-m,\ldots,n-1$ occur at Hessian degree $m$.

\Needspace{24\baselineskip}
\begin{theorem}\label{hessian:global-comparison-characterization}\index{Hessian measure!invariant mass}\index{Comparison principle!Hessian power identities}
The following conditions are equivalent:
\begin{enumerate}
\item $\ddc(\omega^{n-q})=0$ for every $1\leq q\leq m$;
\item $\int_XH_m(u)$ is independent of the smooth admissible
      potential $u$;
\item $\mathcal M_m(t_1,\ldots,t_m)=\int_X\omega^n$ for every
      bounded admissible list;
\item every pair of bounded admissible $u,v$ satisfies
\[
 \int_{\{u<v\}}H_m(v)\leq\int_{\{u<v\}}H_m(u).
\]
\end{enumerate}
Under these conditions, for every $1\leq r\leq m$ and every fixed bounded
admissible list $z_1,\ldots,z_{m-r}$, one also has
\[
 \int_{\{u<v\}}\mathrm d\mu_r(v)
           \leq\int_{\{u<v\}}\mathrm d\mu_r(u).
\]
\end{theorem}
\begin{proof}
The first condition implies the third by
\eqref{eq:hessian-bounded-mass-expansion}. The third implies the
second by taking all inputs equal to a smooth potential.

To prove that the second implies the first, fix any smooth real
$f$. For all sufficiently small real $t$ of either sign,
$\omega+t\ddc f$ is positive, so $tf$ is strictly admissible.
The polynomial
\[
 \int_X(\omega+t\ddc f)^m\wedge\omega^{n-m}
 =\sum_{q=0}^m\binom mq t^q
              \int_X(\ddc f)^q\wedge\omega^{n-q}
\]
is constant on an interval. Consequently every coefficient with
$q\geq1$ vanishes. Polarizing that coefficient by replacing $f$
with $\sum_{a=1}^q s_af_a$ gives
\[
 \int_X\ddc f_1\wedge\cdots\wedge\ddc f_q
                        \wedge\omega^{n-q}=0
\]
for all smooth real $f_a$. Integration by parts in the first
factor, using \cref{lem:ibp}, gives
\[
 \int_Xf_1\ddc f_2\wedge\cdots\wedge\ddc f_q
                         \wedge\ddc(\omega^{n-q})=0.
\]
Since $f_1$ is arbitrary, the smooth top-degree form paired with
it vanishes pointwise. At each point the other complex Hessians
can be prescribed by real quadratic functions times a cutoff.
Complexifying this multilinear identity gives all coordinate
$(q-1,q-1)$-forms, as in the proof of
\cref{thm:mass-invariance-characterization}. They detect every
coefficient of $\ddc(\omega^{n-q})$, so that form vanishes. When
$q=1$, there are no other Hessians and the same conclusion follows
directly from the arbitrariness of $f_1$.

Under the third condition, $\mu_r(w)$ and $\mu_r(u)$ have the same
total mass. \Cref{hessian:comparison-mass-defect} proves the mixed
comparison assertion, and $r=m$ gives the fourth condition.
Finally, suppose the fourth condition holds and fix smooth admissible
$u,v$. Choose $C$ so large that $u<v+C$ everywhere.
Constant invariance gives $H_m(v+C)=H_m(v)$, so comparison gives
$\int_XH_m(v)\leq\int_XH_m(u)$. Exchanging $u,v$ gives equality,
which is the second condition.
\end{proof}

The same varying background that obstructs Monge--Amp\`ere comparison
also obstructs comparison at each intermediate Hessian degree.

\Needspace{20\baselineskip}
\begin{example}\label{hessian:torus-comparison-obstruction}\index{Torus!Hessian comparison obstruction}
Let $X=(\mathbb C/(\mathbb Z+\mathrm i\mathbb Z))^n$, with
$z_j=x_j+\mathrm i y_j$, and put
\[
 \beta_j=\frac{\mathrm i}{2\pi}\,\mathrm dz_j\wedge\mathrm d\bar z_j,
 \qquad b=2+\cos(2\pi x_1),
 \qquad\omega=\beta_1+b\beta_2+\sum_{j=3}^n\beta_j.
\]
For $|t|<\pi^{-2}$ the potential $u_t=t\cos(2\pi x_1)$ is smooth
strictly admissible, and, with $B_n=\int_X\beta_1\wedge\cdots\wedge\beta_n$,
\[
 \int_X H_m(u_t)=n!\left(2-\frac{m\pi^2t}{2n}\right)B_n.
\]
Thus every $-\pi^{-2}<t<0$, together with a constant $C>|t|$,
gives a pair $u=0$, $v=u_t+C$ that violates the unmodified global
comparison inequality.
\end{example}
\begin{proof}
As in \cref{ex:torus-mass},
$\ddc u_t=-\pi^2t\cos(2\pi x_1)\beta_1$.
The coefficient of $\beta_1$ in $\omega+\ddc u_t$ is strictly
positive in the stated range, and all the other coefficients are
positive. Since $(\ddc u_t)^2=0$, expansion gives
\[
 H_m(u_t)=\omega^n+m\ddc u_t\wedge\omega^{n-1}
 =n!b\left(1-\frac mn\pi^2t\cos(2\pi x_1)\right)
                         \beta_1\wedge\cdots\wedge\beta_n.
\]
The period averages of $b$ and $b\cos(2\pi x_1)$ are respectively
$2$ and $1/2$, proving the mass formula. For $t<0$ this mass is
strictly larger than $\int_XH_m(0)=2n!B_n$. Meanwhile $C>|t|$
gives $\{0<u_t+C\}=X$, so the comparison inequality for this pair
would assert the opposite ordering of these masses.
\end{proof}

\subsection{Domination and proper equations}
\label{subsec:hessian-bounded-domination}

The positive background term in a convex rescaling makes a strict
comparison set visible to volume. Under the power identities of
\cref{hessian:global-comparison-characterization}, the global comparison
principle then yields domination for bounded potentials. Canonical
representatives turn the resulting almost everywhere order into
pointwise order.

\Needspace{14\baselineskip}
\begin{theorem}\label{hessian:strict-bounded-domination}\index{Domination principle!bounded Hessian potentials}
Assume $\ddc(\omega^{n-q})=0$ for $1\leq q\leq m$. Let $u,v$
be bounded admissible. If, for some $0\leq c<1$,
\[
 \mathbf1_{\{u<v\}}H_m(u)
                   \leq c\mathbf1_{\{u<v\}}H_m(v),
\]
then $u\geq v$ everywhere. In particular, this conclusion holds
if $H_m(u)(\{u<v\})=0$.
\end{theorem}
\begin{proof}
Choose $B>\|v\|_\infty$ and, for $0<\alpha<1$ with $\alpha^m>c$,
put $v_\alpha=\alpha v-(1-\alpha)B$. Then $v_\alpha\leq v$ and
$E_\alpha=\{u<v_\alpha\}\subset\{u<v\}$.
By \cref{hessian:global-comparison-characterization} and the
assumed restriction inequality,
\[
 \int_{E_\alpha}H_m(v_\alpha)
        \leq\int_{E_\alpha}H_m(u)
        \leq c\int_{E_\alpha}H_m(v).
\]
On the other hand, \cref{hessian:convex-background-term} gives
\[
 (\alpha^m-c)\int_{E_\alpha}H_m(v)
             +(1-\alpha)^m\int_{E_\alpha}\omega^n\leq0.
\]
Both summands are nonnegative, and the coefficient of the volume
is positive. Thus $E_\alpha$ has zero volume. Applying
\cref{hessian:canonical-order} gives $u\geq v_\alpha$ pointwise.
Let $\alpha$ increase to one. The boundedness of $v$ gives
$v_\alpha\to v$ uniformly, so $u\geq v$ everywhere. The final
assertion is the case $c=0$.
\end{proof}

\Needspace{13\baselineskip}
\begin{corollary}\label{hessian:proper-bounded-uniqueness}\index{Hessian equation!bounded proper uniqueness}
Assume $\ddc(\omega^{n-q})=0$ for $1\leq q\leq m$ and fix $a>0$.
If bounded admissible $u,v$ satisfy
\[
 e^{-av}H_m(v)\geq e^{-au}H_m(u),
\]
then $u\geq v$. Consequently, for any fixed finite positive Radon
measure $\nu$, the equation $H_m(u)=e^{au}\nu$ has at most one
bounded admissible solution.
\end{corollary}
\begin{proof}
For $\delta>0$, put $E_\delta=\{u<v-\delta\}$. Multiplying the
assumed measure inequality by the positive bounded Borel function
$e^{au}$ gives on $E_\delta$
\[
 H_m(u)\leq e^{a(u-v)}H_m(v)
                 \leq e^{-a\delta}H_m(v-\delta).
\]
The last identity uses invariance under constants. Since
$e^{-a\delta}<1$, \cref{hessian:strict-bounded-domination} applied
to $u,v-\delta$ gives $u\geq v-\delta$ everywhere. Letting
$\delta\downarrow0$ gives $u\geq v$. Two solutions of the stated
equation have equal weighted measures $e^{-au}H_m(u)=\nu$.
Applying the order assertion in both directions proves uniqueness.
\end{proof}

A strictly increasing exponent also controls changes in logarithmic
data for a fixed Radon right-hand measure.

\Needspace{20\baselineskip}
\begin{proposition}\label{hessian:proper-bounded-stability}\index{Hessian equation!bounded data stability}
Assume $\ddc(\omega^{n-q})=0$ for $1\leq q\leq m$. Let $A<B$ and $f$ be
smooth near $X\times[A,B]$, with $\partial_t f\geq\kappa>0$
there, and let $\nu$ be a finite positive Radon measure.
Suppose bounded admissible $u,v$ have values in $[A,B]$ and satisfy
\[
 H_m(u)=e^{f(x,u)+h_1}\nu,\qquad
 H_m(v)=e^{f(x,v)+h_2}\nu,
\]
where $h_1,h_2$ are bounded Borel functions. Then
\[
 \|u-v\|_\infty
            \leq\kappa^{-1}\|h_1-h_2\|_{L^\infty(\nu)}.
\]
\end{proposition}
\begin{proof}
Put $K=\|h_1-h_2\|_{L^\infty(\nu)}$ and fix $s>K/\kappa$.
On $E_s=\{v+s<u\}$, the actual values $u,v$ lie in $[A,B]$,
so $f(x,u)-f(x,v)\geq\kappa s$. The equations and the essential
bound on $h_2-h_1$ therefore give, as restricted measures,
\[
 \mathbf1_{E_s}H_m(v+s)
 =\mathbf1_{E_s}H_m(v)
 \leq e^{K-\kappa s}\mathbf1_{E_s}H_m(u).
\]
The constant is smaller than one. Applying
\cref{hessian:strict-bounded-domination} to $v+s,u$ gives
$u\leq v+s$ pointwise. Let $s\downarrow K/\kappa$ and exchange
$u,v$ to obtain the norm bound. The exponent is evaluated only
at $u,v$; the shifted comparison potential need not take its
values in $[A,B]$.
\end{proof}

\Needspace{14\baselineskip}
\begin{corollary}\label{hessian:proper-bounded-smoothness}\index{Hessian equation!bounded proper smoothness}
Assume $\ddc(\omega^{n-q})=0$ for $1\leq q\leq m$.
Under the hypotheses of \cref{hessian:proper-variable-existence},
every bounded admissible solution with values in $[A,B]$ is its
smooth solution. In particular, for every $a>0$ and smooth real
$H$, the equation $H_m(u)=e^{au+H}\omega^n$ has exactly one bounded
admissible solution, and that solution is smooth and strictly
admissible.
\end{corollary}
\begin{proof}
For a continuous admissible potential, smooth decreasing admissible
approximations converge uniformly on compact $X$ by Dini's theorem.
Thus \cref{hessian:joint-decreasing-products,hessian:continuous-products}
define the same Hessian measure on the continuous class.
The smooth solution of \cref{hessian:proper-variable-existence}
therefore solves the bounded equation as well. Apply
\cref{hessian:proper-bounded-stability} with $\nu=\omega^n$ and
$h_1=h_2=0$ to identify any other bounded solution in $[A,B]$
with it. For the final assertion, smooth existence is already
\cref{hessian:smooth-proper-existence}, and
\cref{hessian:proper-bounded-uniqueness} with
$\nu=e^H\omega^n$ gives uniqueness among all bounded solutions.
\end{proof}
